% Benchmarking and evaluation: forecasting evaluation protocols (one-page version).
% Long version: evaluation_protocols_long.tex. References: evaluation_protocols.bib.
% Numbers: aeon 8585c7bf, results/aeon-8585c7bf/{m4,aaltd} in the horizon repo.
% Table labels (tab:m4, tab:aaltd) are placeholders.

\subsection{Forecasting evaluation protocols}
\label{sec:protocols}

There is no standard protocol for evaluating forecasters. A protocol fixes the
forecast origins, the points at which the history given to a model ends and
its forecast begins (a single origin per series, which is then the train/test
boundary, or an origin rolled forward through the test period), whether models
are refitted at each origin, the horizon, how multi-step forecasts are made
(recursively, feeding one-step predictions back in; directly, one model per
step; or one model for the whole horizon~\citep{taieb2012review}), whether
models are local, global or pretrained, and how errors are scaled and
aggregated~\citep{tashman2000outofsample,hewamalage2023forecast}. Each choice
can change the ranking. aeon makes the strategy explicit
(\texttt{iterative\_forecast}, \texttt{direct\_forecast},
\texttt{series\_to\_series\_forecast}), so one model can be run under different
protocols.

M4~\citep{makridakis2020m4} holds out the last $h$ points of each of 100{,}000
series ($h$ from 6 for yearly to 48 for hourly) and forecasts once from that
single origin. Most series are short (yearly series need only 13 observations)
and 95\% are yearly, quarterly or monthly. Accuracy is OWA, the mean of sMAPE
and MASE relative to Naive2, pooled over all series, so weekly and hourly
series (under 1\%) carry almost no weight; the benchmarks, including the ML
ones, were run on series deseasonalised by a 90\% autocorrelation test. The
series are not aligned in time, so global models can leak information across
series~\citep{ingel2020correlated}. The Monash
repository~\citep{godahewa2021monash} keeps this single fixed origin for a
wider set of longer, higher-frequency datasets, with competition horizons where
they exist and frequency-based ones otherwise.

Deep learning papers usually evaluate on M4 differently.
TimesNet~\citep{wu2023timesnet}, through the Time-Series-Library, follows
N-BEATS~\citep{oreshkin2020nbeats}: one global model per frequency is trained
on windows sampled from the end of all the series, forecasts the whole horizon
directly from only the last $2h$ observations, and is scored with weekly,
daily and hourly pooled into one ``Others'' group. Its validation set is the
M4 test set, so early stopping and checkpoint selection use the test
values.\footnote{Time-Series-Library, \texttt{Dataset\_M4} loads the test split
for any flag other than \texttt{train}, and \texttt{vali} scores forecasts
against it; present from the first commit (\texttt{e08ddc0}, February 2023),
the code release for TimesNet.} The test points and the
OWA formula match M4, but the information each model sees does not:
statistical models use the full history of one series, the deep models a short
window of every series.

GIFT-Eval~\citep{aksu2024gifteval} mixes protocols within one score. The last
10\% of each series is split into up to 20 non-overlapping windows, giving a
rolling origin (M4 keeps its single origin), and horizons are scaled by 1, 10
and 15 (up to 720 hourly steps). Statistical models are refitted per window on
at most the last 1{,}000 points, deep models are trained once per dataset, and
foundation models forecast zero-shot, the decoder-only ones (Chronos, TimesFM)
recursively, which the authors link to their decline at long horizons.

We evaluate aeon's forecasters under two contrasting protocols with the same
code: M4 exactly as in the competition (local models, one origin, recursive
one-step machine learning models), and 100 long series (up to 10{,}000 points,
mostly from the Monash repository, including 57 long M4 series) with a rolling
origin, refitting at each of the last 30 points and forecasting one step ahead.
On M4 the statistical forecasters win (SCUM OWA 0.845) and every machine
learning forecaster is worse than Naive2 (DrCIF 1.450); on the long series
DrCIF is best (MAE relative to naive 0.781, SCUM 0.843)
(Tables~\ref{tab:m4},~\ref{tab:aaltd}). Neither ranking is wrong; each answers
a different question.
